\section{A low-degree estimate on biased products}\label{sec:fourier}

We bound the mass of low-degree Fourier coefficients for Boolean
functions whose individual influences are small compared with their
total influence. The bound is linear in total influence, so it can be
averaged over the conditioned coordinates in Section~\ref{sec:graph}.

Let $J$ be a finite coordinate set, let $0<p<1$, and put
$\sigma=\sqrt{p(1-p)}$. Expectations and norms in this section are with
respect to the product Bernoulli-$p$ measure on the indicated coordinates;
in particular, $\|g\|_q=(\E|g|^q)^{1/q}$. Define
\[
  \chi_i(x)=\frac{x_i-p}{\sigma},\qquad
  \chi_S(x)=\prod_{i\in S}\chi_i(x),\qquad \chi_\varnothing=1.
\]
Each $\chi_i$ has mean zero and second moment one. Independence therefore
makes the $2^{|J|}$ functions $\chi_S$, $S\subseteq J$, orthonormal.
They form a basis of the $2^{|J|}$-dimensional space of real functions on
$\{0,1\}^J$. Thus, with $\widehat g(S)=\E[g\chi_S]$,
\begin{equation}\label{eq:parseval}
  g=\sum_{S\subseteq J}\widehat g(S)\chi_S,\qquad
  \E g^2=\sum_{S\subseteq J}\widehat g(S)^2,\qquad
  \Var(g)=\sum_{\varnothing\ne S\subseteq J}\widehat g(S)^2.
\end{equation}

For $i\in J$, write
$\Delta_i g(x)=g(x_{J\setminus\{i\}},1)-g(x_{J\setminus\{i\}},0)$,
a function on the remaining coordinates. Integrating just coordinate $i$
gives
\[
  \E_i[g\chi_i]
  =\frac{p(1-p)}{\sigma}
    \bigl(g(x_{J\setminus\{i\}},1)-g(x_{J\setminus\{i\}},0)\bigr)
  =\sigma\Delta_i g.
\]
Consequently, for $A\subseteq J\setminus\{i\}$,
$\widehat g(A\cup\{i\})=\sigma\widehat{\Delta_i g}(A)$,
where the latter Fourier coefficient is taken on $J\setminus\{i\}$.
For Boolean $h:\{0,1\}^J\to\{0,1\}$, set
\[
  I_i(h)=\Pr(\Delta_i h\ne0),\qquad I(h)=\sum_{i\in J}I_i(h).
\]
Since $\Delta_i h\in\{-1,0,1\}$, Parseval on the remaining coordinates
and then summation over $i$ give
\begin{equation}\label{eq:dirichlet}
  \sum_{S\ni i}\widehat h(S)^2=\sigma^2 I_i(h),\qquad
  \sum_{S\subseteq J}|S|\widehat h(S)^2=\sigma^2 I(h).
\end{equation}
These identities do not require monotonicity.

The next estimate follows the hypercontractive influence method of
Kahn, Kalai, and Linial~\cite{KKL88}. We prove the biased two-point
inequality directly and use the standard product argument of
Bonami~\cite[Chapter~III, Lemma~1]{Bonami1970}, keeping the dependence
on the edge probability explicit.

\begin{samepage}
\begin{lemma}\label{lem:low-degree}
Let $m\ge2$ be real and let $h:\{0,1\}^J\to\{0,1\}$ satisfy
$I_i(h)\le I(h)/m$ for every $i\in J$. If
\[
  k=\frac{\sigma\log m}{64},
\]
then
\begin{equation}\label{eq:low-degree}
  \sum_{1\le |S|\le k}\widehat h(S)^2\le m^{-1/4}I(h).
\end{equation}
Every degree cutoff is interpreted literally for integer degrees $|S|$;
in particular, the sum is empty when $k<1$.
\end{lemma}
\end{samepage}

\begin{proof}
We first prove a product norm bound. Put $\rho=\sigma/4$, so
$0<\rho\le1/8$, and define the linear operator $T_\rho$ by
$T_\rho\chi_S=\rho^{|S|}\chi_S$. For one coordinate, the bound
$|\chi_i|\le\sigma^{-1}$ and the second moment identity imply
\[
  |\E\chi_i^3|\le\sigma^{-1},\qquad
  \E\chi_i^4\le\sigma^{-2}.
\]
For real $a,b$, expansion and
$2|ab^3|\le a^2b^2+b^4$ yield
\begin{align*}
  \E(a+\rho b\chi_i)^4
  &\le a^4+6\rho^2a^2b^2
       +\frac{4\rho^3}{\sigma}|ab^3|
       +\frac{\rho^4}{\sigma^2}b^4\\
  &\le a^4+\frac{13\sigma^2}{32}a^2b^2
       +\frac{9\sigma^2}{256}b^4
   \le(a^2+b^2)^2,
\end{align*}
where the last inequality uses $\sigma^2\le1/4$.

To tensorize this inequality, induct on $|J|$, the empty product being
immediate. Write $g=a+b\chi_i$ with $a,b$ functions of the other
coordinates, and set $A=T_\rho a$, $B=T_\rho b$, using the operator on
those coordinates. Apply the scalar inequality pointwise and then the
triangle inequality in $L^2$ to obtain
\[
  \|T_\rho g\|_4^2
  \le \|A^2+B^2\|_2
  \le \|A\|_4^2+\|B\|_4^2
  \le \|a\|_2^2+\|b\|_2^2
  =\|g\|_2^2.
\]
The penultimate inequality is the induction hypothesis. We have proved
\begin{equation}\label{eq:noise-norm}
  \|T_\rho g\|_4\le\|g\|_2
\end{equation}
on every finite product.

For any real $d\ge0$, let $P_{\le d}$ be the orthogonal projection onto
the span of the characters of degree at most $d$. If $u=P_{\le d}u$,
the operator $T_\rho$ is invertible and Parseval gives
\[
  \|T_\rho^{-1}u\|_2^2
  =\sum_{|S|\le d}\rho^{-2|S|}\widehat u(S)^2
  \le\rho^{-2d}\|u\|_2^2.
\]
Applying \eqref{eq:noise-norm} to $T_\rho^{-1}u$ proves
$\|u\|_4\le\rho^{-d}\|u\|_2$. Duality in this finite-dimensional
subspace, followed by H\"older's inequality, now gives for every real
function $v$
\begin{equation}\label{eq:projection-norm}
  \|P_{\le d}v\|_2
  =\sup_{\substack{u=P_{\le d}u\\\|u\|_2=1}}|\E[vu]|
  \le\rho^{-d}\|v\|_{4/3}.
\end{equation}

If $k<1$, \eqref{eq:low-degree} is immediate, so assume $k\ge1$.
Since $|\Delta_i h|$ is the indicator of a pivotal event,
$\|\Delta_i h\|_{4/3}=I_i(h)^{3/4}$. Every nonempty Fourier set is
counted at least once when summing over its coordinates. The derivative
coefficient identity and \eqref{eq:projection-norm}, applied on
$J\setminus\{i\}$ with $d=k$, therefore show
\begin{align*}
  \sum_{1\le |S|\le k}\widehat h(S)^2
  &\le\sum_{i\in J}\sum_{\substack{S\ni i\\|S|\le k}}
       \widehat h(S)^2\\
  &=\sigma^2\sum_{i\in J}
       \sum_{\substack{A\subseteq J\setminus\{i\}\\|A|\le k-1}}
       \widehat{\Delta_i h}(A)^2\\
  &\le\sigma^2\rho^{-2k}\sum_{i\in J}I_i(h)^{3/2}.
\end{align*}
The function $s\mapsto s\log(4/s)$ is increasing on $(0,1/2]$, since
its derivative is $\log(4/s)-1>0$. Hence
$\sigma\log(4/\sigma)\le\tfrac12\log8<2$, and
\[
  \rho^{-2k}
  =m^{\sigma\log(4/\sigma)/32}\le m^{1/16}.
\]
If $I(h)\le m^{1/4}$, the hypothesis implies
\[
  \sum_i I_i(h)^{3/2}
  \le\sqrt{\frac{I(h)}m}\,I(h)
  \le m^{-3/8}I(h).
\]
The preceding estimates bound the desired sum by
$\sigma^2m^{-5/16}I(h)\le m^{-1/4}I(h)$, including when $I(h)=0$.
If $I(h)>m^{1/4}$, Parseval and Booleanity instead give
\[
  \sum_{1\le |S|\le k}\widehat h(S)^2
  \le\E h^2\le1<m^{-1/4}I(h).
\]
This proves the estimate for every $0<p<1$, without a lower bound on
$\sigma$.
\end{proof}
